(1) 在泵浦不耗尽假设下直接积分：
\[
A_2(L)=i\kappa A_1^2\int_0^Le^{i\Delta kz}dz
=i\kappa A_1^2Le^{i\Delta kL/2}
\operatorname{sinc}(\Delta kL/2),
\]
其中 \(\operatorname{sinc}x=\sin x/x\) 且在 x=0 取连续极限1。

(2) 输出倍频光强与 \(|A_2|^2\) 成正比，所以
\[
|A_2(L)|^2=\kappa^2|A_1|^4L^2
\operatorname{sinc}^2(\Delta kL/2).
\]
以相同长度下 Δk=0 的强度为基准，归一化转换效率为 \(\eta=\operatorname{sinc}^2(\Delta kL/2)\le1\)。

(3) 当相位完全匹配 Δk=0，各微小晶体切片所产生的二次谐波相位一致，复振幅随 L 线性增长为 \(A_2=i\kappa A_1^2L\)，强度则随 \(L^2\) 增长；不能把相干振幅误按强度进行简单求和。

(4) 通常把相对相位累积 π 所对应距离定义为相干长度
\[
L_{\rm coh}=\frac{\pi}{|\Delta k|}.
\]
若晶长为 \(2L_{\rm coh}\)，输出振幅中的 sinc 因子为零，表明后半部分产生的倍频场可与前半部分相消。该完美相消结论严格属于本题的无耗散恒泵浦、均匀晶体模型；当实际倍频强烈导致泵浦耗尽或材料色散、吸收显著时必须求解完整耦合波方程。
